\documentclass [english,memo]{xarticle}
%
\usepackage[english]{babel}
\usepackage{cpbfonts}
\renewcommand{\ttdefault}{lmtt}
\usepackage{url}
\usepackage{cpbextra}
%
\makeatletter
\def\section@font{\head@style\normalfont\sffamily\large\bfseries}
\makeatother
%
\usepackage{needspace}
%
\usepackage{xspace}
\usepackage{shortvrb}
\usepackage{amsmath}
\usepackage[psamsfonts]{amsfonts}

\MakeShortVerb{\|}
%
\newcommand*{\IsisTool}{\textsf{IsisTool}\xspace}
\newcommand*{\Deckisis}{\texttt{Deckisis}\xspace}
%
\newcommand*{\FtnWord}[1]{\texttt{#1}}
\newcommand{\Rbb}{\mathbb{R}}
%
\newenvironment{leftindent}%
{\begin{list}{}{\setlength{\leftmargin}{1cm}}\item[]}%
{\end{list}}

%
%
%

\begin{document}

\To{Isis users}
\author{Rob van Harrevelt}
\date{\today}
\Subject{The SVD analysis for rank loss in matrices}
%
\maketitle
%


\section{Introduction}

\section{The SVD}

The SVD of a matrix $A \in \Rbb^{m \times n}$ is given by (see \cite{matcomp})
\begin{equation}
A = U \: \Sigma \: V^T
\end{equation}
where $U \in \Rbb^{m \times m}$ and  $U \in \Rbb^{n \times n}$ are
orthogonal matrices and the matrix $\Sigma \in \Rbb^{m \times n}$
is given by
\begin{equation}
\Sigma_{ij} = \begin{cases}
               \sigma_i & \text{if $i = j$} \\
	       0 & \text{if $i \neq j$}
              \end{cases}
\end{equation}
The $\sigma_i$'s, called the singular values, are ordered such that
\begin{equation}
\sigma_1 \ge \sigma_2 \ge ... \ge \sigma_p \ge 0  \qquad p = \text{min}\{m, n\}
\end{equation}

The column vectors $u_i$ and $v_j$ are called the left and right singular
vectors respectively. These vectors obey the relation
(see \cite{matcomp})
\begin{equation}\label{eq:singular_vectors}
\left.
\begin{array}{llcl}
A   & v_i & = & \sigma_i u_i \\
A^T & u_i & = & \sigma_i v_i
\end{array}
\right\}
i = \text{min}\{m, n\}
\end{equation}

\section{Rank loss and SVD analysis}

The rank $r$ of the matrix $A$ is equal to the number of non-zero singular values.
Thus $r \leq \text{min}\{m, n\}$.
If case $r < \text{min}\{m, n\}$ is called rank loss.

If $\sigma_k = 0$, then from Eq. \ref{eq:singular_vectors} it follows that
\begin{equation}\label{eq:null_vectors}
\begin{array}{llcl}
A   &v_k & = & 0 \\
A^T & u_k & = & 0
\end{array}
\end{equation}

Thus the vector $v_k$ is the linear combination of columns of
matrix $A$ that yields 0, and the vector $u_k$ is the linear combination
of rows of $A$ that yields zero. The vector $v_k$ and $u_i$ thus
give information about linearly dependent rows and columns in matrix $A$.

\bibliography{svd_analysis}
\bibliographystyle{cpblit}

\end{document}
